% Ampliación sustantiva de la completitud coinductiva Hodge.
% Módulo destinado al capítulo Hodge de la Parte IV.

\section{Coinductive unfolding of Hodge completeness}
\label{sec:amp-hodge-completitud-desplegada}

The completeness equality
\eqref{eq:hodge-suc-completitud} can be unfolded without starting from an algebraic
class or introducing into the source the conclusion one wishes to publish. The
object being completed is the system of enriched histories; the Hodge
class intervenes only as compatible evaluation data in its
finite truncations.

\subsection{Finite observation systems and compatible fibers}
\label{subsec:amp-hodge-sistemas-finitos}

For each depth \(N\geq0\), let \(J_N(X,p)\) be the finite category of
incidence charts, oriented routes and memory terminals of depth at
most \(N\). We write
\begin{equation}\label{eq:amp-hodge-surv-n}
 \mathscr S_N(X,p):=\operatorname{Surv}_{J_N}(X,p),
 \qquad
 \rho_{N+1,N}:\mathscr S_{N+1}(X,p)\longrightarrow\mathscr S_N(X,p)
\end{equation}
for the rationalized space of surviving states and its restriction. The
finite readers take values in the module of cylinder observations
\(\mathscr H_N^p(X)\):
\begin{equation}\label{eq:amp-hodge-evaluacion-finita}
 e_N^{\mathrm{Hdg}}:\mathscr S_N(X,p)\longrightarrow
 \mathscr H_N^p(X),
 \qquad
 q_{N+1,N}:\mathscr H_{N+1}^p(X)\longrightarrow\mathscr H_N^p(X).
\end{equation}
Here \(\mathscr H_N^p(X)\) records only the evaluations on the incidences
present in \(J_N\), but preserves their rational coefficients and
orientation. If
\(q_N:\operatorname{Hdg}^p(X)_{\mathbb Q}\to\mathscr H_N^p(X)\) is the
restriction to the same finite set of observations, we have
\begin{equation}\label{eq:amp-hodge-cuadrado-restriccion}
 e_N^{\mathrm{Hdg}}\rho_{N+1,N}
 =q_{N+1,N}e_{N+1}^{\mathrm{Hdg}},
 \qquad
 q_{N+1,N}q_{N+1}=q_N.
\end{equation}

For \(\alpha\in\operatorname{Hdg}^p(X)_{\mathbb Q}\), define the compatible finite
fiber
\begin{equation}\label{eq:amp-hodge-fibra-alpha}
 \mathscr F_N(\alpha):=
 \left\{s_N\in\mathscr S_N(X,p):
 e_N^{\mathrm{Hdg}}(s_N)=q_N(\alpha)\right\}.
\end{equation}
The word \emph{finite} refers to the depth and the incidence
diagram, not to a loss of rational coefficients or to a selection
of a representative from \(X\) alone.

\begin{proposition}[Finite extension with memory]
\label{prop:amp-hodge-extension-finita}
For every \(\alpha\) as in \eqref{eq:amp-hodge-fibra-alpha}, the fibers
\(\mathscr F_N(\alpha)\) are nonempty and the restrictions induce surjective
maps
\begin{equation}\label{eq:amp-hodge-mittag-leffler}
 \rho_{N+1,N}:\mathscr F_{N+1}(\alpha)\twoheadrightarrow
 \mathscr F_N(\alpha).
\end{equation}
The extension preserves sheet, TRIT sign, quotient, carry, support, boundary and
memory terminal.
\end{proposition}

\begin{proof}
The assertion is the Hodge specialization of the extension theorem for the atlas
constructed in Part II. Its mechanism should be made visible. Let
\(s_N\in\mathscr F_N(\alpha)\) and choose a prolongation of charts
\(\widetilde s_{N+1}\) preserving the history of \(s_N\). Its discrepancy at
the new level is
\begin{equation}\label{eq:amp-hodge-discrepancia}
 \delta_{N+1}:=q_{N+1}(\alpha)
 -e_{N+1}^{\mathrm{Hdg}}(\widetilde s_{N+1}),
 \qquad q_{N+1,N}(\delta_{N+1})=0,
\end{equation}
where the last equality comes from
\eqref{eq:amp-hodge-cuadrado-restriccion}. The
APP--Mayer--Vietoris cells publish exactly that relative kernel through the
cones of the morphisms \(\kappa_{a,b}\); their defect
\((a-1)(b-1)\) is lifted with the residue and carry of
\eqref{eq:hodge-suc-celda-carry}. By the finite extension theorem there is a
relative correction \(\eta_{N+1}(\delta_{N+1})\), supported on the new
charts, such that
\begin{equation}\label{eq:amp-hodge-extension-operativa}
 \rho_{N+1,N}\eta_{N+1}(\delta_{N+1})=0,
 \qquad
 e_{N+1}^{\mathrm{Hdg}}\eta_{N+1}(\delta_{N+1})=\delta_{N+1}.
\end{equation}
Therefore
\begin{equation}\label{eq:amp-hodge-ext}
 \operatorname{Ext}_{N+1}(s_N,\alpha):=
 \widetilde s_{N+1}+\eta_{N+1}(\delta_{N+1})
 \in\mathscr F_{N+1}(\alpha)
\end{equation}
restricts to \(s_N\). The orientation
\(\kappa_{b,a}\tau=-P\kappa_{a,b}\) fixes the sign, the cocycle law
\eqref{eq:hodge-suc-carry} fixes the composition and
\(\operatorname{Cone}(I-U_{\Gamma_9})\) preserves the terminal. The initial level
is obtained from the based seed atlas. Induction proves nonemptiness and
surjectivity at all levels.
\end{proof}

\subsection{Passage to the limit and order of construction}
\label{subsec:amp-hodge-paso-limite}

\begin{theorem}[Unfolded coinductive completeness]
\label{thm:amp-hodge-completitud-coinductiva}
With the data already constructed in Part II---the systems
\eqref{eq:amp-hodge-surv-n}--\eqref{eq:amp-hodge-evaluacion-finita}, the
compatibility \eqref{eq:amp-hodge-cuadrado-restriccion}, the finite extension
of proposition \ref{prop:amp-hodge-extension-finita}, the separation of
cylinder evaluations and preservation of the terminal---for each
\(\alpha\in\operatorname{Hdg}^p(X)_{\mathbb Q}\), there is an enriched
history
\begin{equation}\label{eq:amp-hodge-xi-limite}
 \xi_\alpha=(s_0,s_1,s_2,\ldots)
 \in\varprojlim_N\mathscr F_N(\alpha)
 \subset X_\chi(X,p)
\end{equation}
which satisfies
\begin{equation}\label{eq:amp-hodge-evaluacion-alpha}
 \operatorname{ev}^{\mathrm{Hdg}}_{\infty,X,p}(\xi_\alpha)=\alpha.
\end{equation}
In particular, completeness is obtained before applying the perfect
realizer and before forming any Chow class.
\end{theorem}

\begin{proof}
Choose \(s_0\in\mathscr F_0(\alpha)\). Once \(s_N\) is constructed, the
surjectivity \eqref{eq:amp-hodge-mittag-leffler} allows us to choose
\(s_{N+1}\in\mathscr F_{N+1}(\alpha)\) with
\(\rho_{N+1,N}s_{N+1}=s_N\). Thus one obtains
\eqref{eq:amp-hodge-xi-limite}. Equivalently, the system of fibers
satisfies the Mittag--Leffler condition and produces no
\(\varprojlim^1\) obstruction term. By the definition of each fiber,
\begin{equation}\label{eq:amp-hodge-evaluaciones-xi}
 e_N^{\mathrm{Hdg}}(s_N)=q_N(\alpha)
 \qquad(N\geq0).
\end{equation}
The separation of cylinder observations identifies the unique element
of the limit of the \(q_N(\alpha)\) with \(\alpha\), which proves
\eqref{eq:amp-hodge-evaluacion-alpha}. The terminal cone prevents replacing the
sequence \((s_N)_N\) with the repetition of a visible root; therefore
\(\xi_\alpha\) preserves the rigidification used in passing to the limit.
\end{proof}

The finite publication is arranged in commutative squares
\begin{equation}\label{eq:amp-hodge-cuadrado-finito}
 \begin{array}{ccc}
 \mathscr S_N(X,p)&\xrightarrow{\ R_{\mathrm{Hdg},N}\ }&
 K_0(\mathbf{Perf}^{\geq p}(X))\otimes\mathbb Q\\[1mm]
 \big\downarrow\scriptstyle e_N^{\mathrm{Hdg}}&&
 \big\downarrow\scriptstyle
   \operatorname{cl}_{X,p}\circ\operatorname{ch}^{\mathrm{loc}}_p\\[1mm]
 \mathscr H_N^p(X)&\xrightarrow{\ \operatorname{pub}_N\ }&
 H^{2p}(X,\mathbb Q).
 \end{array}
\end{equation}
The naturality of these squares with respect to \(\rho_{N+1,N}\) gives, in the
limit, the identity \eqref{eq:hodge-suc-identidad}. The causal order is,
therefore,
\begin{equation}\label{eq:amp-hodge-orden-causal}
 \alpha\longmapsto(q_N\alpha)_N
 \longmapsto\xi_\alpha
 \longmapsto R_{\mathrm{Hdg}}(\xi_\alpha)
 \longmapsto
 \beta_\alpha:=\operatorname{ch}^{\mathrm{loc}}_p
   R_{\mathrm{Hdg}}(\xi_\alpha).
\end{equation}
Only in the last step does \(\beta_\alpha\) appear. Applying the limit square
to \eqref{eq:amp-hodge-evaluacion-alpha} yields
\begin{equation}\label{eq:amp-hodge-publicacion-final}
 \operatorname{cl}_{X,p}(\beta_\alpha)=
 \operatorname{ev}^{\mathrm{Hdg}}_{\infty,X,p}(\xi_\alpha)=\alpha.
\end{equation}

\subsection{Local check on the quadric
\texorpdfstring{\(Q^4\)}{Q4}}
\label{subsec:amp-hodge-q4}

The split quadric
\begin{equation}\label{eq:amp-hodge-q4}
 Q^4=\{x_0x_3+x_1x_4+x_2x_5=0\}\subset\mathbb P^5_{\mathbb C}
\end{equation}
has two based families of maximal planes, with representatives
\(\Pi_+\) and \(\Pi_-\). If
\(a=\operatorname{cl}[\Pi_+]\) and
\(b=\operatorname{cl}[\Pi_-]\), then
\begin{equation}\label{eq:amp-hodge-q4-bases}
 \operatorname{CH}^2(Q^4)=\mathbb Za\oplus\mathbb Zb,
 \qquad
 H^4(Q^4,\mathbb Z(2))\cap H^{2,2}(Q^4)
 =\mathbb Za\oplus\mathbb Zb.
\end{equation}
The based map
\begin{equation}\label{eq:amp-hodge-q4-beta}
 \beta_{Q^4,2}(ra+sb):=r[\Pi_+]+s[\Pi_-]
\end{equation}
satisfies exactly
\begin{equation}\label{eq:amp-hodge-q4-beta-identidades}
 \partial_{\mathrm{mot}}\beta_{Q^4,2}=0,
 \qquad
 \operatorname{cl}_{Q^4}\beta_{Q^4,2}=I.
\end{equation}

For the six ports, let
\begin{equation}\label{eq:amp-hodge-q4-matriz}
 A=\begin{pmatrix}
 2&2&2&1&2&1\\
 2&1&2&2&1&1\\
 1&0&0&1&0&2\\
 0&0&1&1&1&0\\
 2&2&2&0&2&1\\
 2&1&2&1&0&0
 \end{pmatrix},
 \qquad \det A=1,
 \qquad \mathbb A=A\otimes I.
\end{equation}
The division \(x_k\mathbb A=u_k+3x_{k+1}\) and the reduction of
\eqref{eq:amp-hodge-q4-beta} define, gate by gate,
\begin{equation}\label{eq:amp-hodge-q4-uk}
 U_{k,j}:=\overline\beta_{Q^4,2}(u_{k,j}),
 \qquad
 D_{\mathrm{mot}}U_k=0,
 \qquad
 H^0(R_{\mathrm{mot}}\otimes\mathbb F_3)(U_k)=u_k.
\end{equation}
Iterating nine times yields the telescoping identity
\begin{equation}\label{eq:amp-hodge-q4-telescopia}
 x_0\mathbb A^9
 =R_{\mathrm{mot}}\!\left(
   \sum_{k=0}^{8}3^kU_k\mathbb A^{8-k}\right)+3^9x_9.
\end{equation}
Equations \eqref{eq:amp-hodge-q4-beta-identidades}--
\eqref{eq:amp-hodge-q4-telescopia} verify local closure, port compatibility
and nonadic memory in a rank-two model. They constitute a
specialized validator of the extension mechanism; general completeness
comes from the coinductive system of the preceding theorems, not from an
extrapolation from \(Q^4\).

\subsection{Premises and refutation criteria}
\label{subsec:amp-hodge-falsadores}

The proof uses exactly five data: the finite atlas of based
charts; the relative extension of proposition
\ref{prop:amp-hodge-extension-finita}; the restriction compatibility
\eqref{eq:amp-hodge-cuadrado-restriccion}; the separation of cylinder
evaluations; and the commutativity of the publication
\eqref{eq:amp-hodge-cuadrado-finito}. Each has a locatable refutation
criterion:
\begin{enumerate}
 \item an empty fiber \(\mathscr F_N(\alpha)\) would refute finite coverage;
 \item failure of \eqref{eq:amp-hodge-mittag-leffler} would prevent constructing the
       compatible history;
 \item two distinct classes with all truncations equal would refute the
       separation of the reader;
 \item a noncommutative finite square would refute
       \eqref{eq:amp-hodge-publicacion-final};
 \item replacing the terminal with a memoryless root would invalidate the
       compatibility of the sequence, even if it preserved its visible value.
\end{enumerate}
These criteria target specific arrows of the argument. The obstruction to a
finite selector depending only on \(X\) affects none of them, because
the construction preserves the based antecedent \(\xi_\alpha\).

\begin{theorem}[Generation of the algebraic class]
\label{thm:hodge-suc-generacion}
For every \(\alpha\in\operatorname{Hdg}^{p}(X)_{\mathbb Q}\) there exists
\(\xi_{\alpha}\in X_{\chi}(X,p)\) such that, on defining
\begin{equation}\label{eq:hodge-suc-beta}
 \beta_{\alpha}:=
 \operatorname{ch}^{\mathrm{loc}}_{p}
 \bigl(R_{\mathrm{Hdg}}(\xi_{\alpha})\bigr)
 \in\operatorname{CH}^{p}(X)_{\mathbb Q},
\end{equation}
we have
\begin{equation}\label{eq:hodge-suc-cl-beta}
 \operatorname{cl}_{X,p}(\beta_{\alpha})=\alpha.
\end{equation}
\end{theorem}

\begin{proof}
By the completeness theorem
\eqref{eq:hodge-suc-completitud}, every \(\alpha\) has an antecedent
\(\xi_{\alpha}\) before applying the perfect reader. One then forms
\(\beta_{\alpha}\) through \eqref{eq:hodge-suc-beta}. The identity
\eqref{eq:hodge-suc-identidad} gives
\[
 \operatorname{cl}_{X,p}(\beta_{\alpha})
 =\operatorname{ev}^{\mathrm{Hdg}}_{\infty,X,p}(\xi_{\alpha})
 =\alpha.
\]
The rigidification lies in \(\xi_{\alpha}\), where base,
orientation, boundary and route are preserved. Therefore, choosing an antecedent does not
become a natural selector depending only on \(X\), nor has the
sought surjectivity been introduced as a premise of the map
\(R_{\mathrm{Hdg}}\).
\end{proof}

\subsection*{Hypotheses and scope}

Theorem \ref{thm:hodge-suc-generacion} uses exactly: the smooth projective
variety \(X\) and the integer \(p\); the complete atlas
\(X_{\chi}(X,p)\) and its completeness theorem
\eqref{eq:hodge-suc-completitud}; the orientations and supports of its charts; the
Tor-independence of the local intersections; the localized Chern character and
the class map defined with the same convention; and preservation of the
terminal \eqref{eq:hodge-suc-terminal}. With these explicit initial
data, \eqref{eq:hodge-suc-cl-beta} is an exact result for the entire
domain \eqref{eq:hodge-suc-dominio}.

The double publication and preservation of history are achieved through the
cell \(\kappa_{a,b}\), its cone, totalization by coend and the map
\(R_{\mathrm{Hdg}}\) with values in \(\operatorname{Perf}\). The identity
\eqref{eq:hodge-suc-identidad} and the generation
\eqref{eq:hodge-suc-cl-beta} are the results of the specialization; their
focal verifiers check support, rank, sign, carry, naturality,
terminal and commutativity of the publication.

\subsection{Final mathematical recognition of Hodge}

The final translation of \eqref{eq:hodge-suc-cl-beta} is
\begin{equation}\label{eq:hodge-suc-reconocimiento-final}
 \operatorname{cl}_{X,p}:
 \operatorname{CH}^{p}(X)_{\mathbb Q}
 \twoheadrightarrow
 \operatorname{Hdg}^{p}(X)_{\mathbb Q}
\end{equation}
for every smooth complex projective variety \(X\) and every \(p\geq0\). This
is exactly the conventional rational formulation of the Hodge
statement. Here it appears as recognition of the chain
\(X_{\chi}\to\operatorname{Perf}\to\operatorname{CH}^{p}\to
\operatorname{Hdg}^{p}\), not as a condition used to construct it.

% Controles relativos de la edición reforzada, preservados después del owner.
\subsection{Relative auxiliary control: coend, kernels and based lifting}
\label{subsec:amp-hodge-presentacion-relativa}

\noindent\textbf{Status: NON\_GOVERNING\_CONTROL.}
The proprietary construction already proved is the APP--Mayer--Vietoris extension
with memory of \cref{prop:amp-hodge-extension-finita}, followed by the
Mittag--Leffler limit of
\cref{thm:amp-hodge-completitud-coinductiva}. The following kernels, cokernels and
coends are an additional relative presentation and local
falsifiers: they do not select \(\xi_\alpha\), are not premises of the proprietary
construction and cannot reopen
\eqref{eq:amp-hodge-evaluacion-alpha}.

To audit redundantly an extension already obtained by proposition
\ref{prop:amp-hodge-extension-finita}, this chart defines a relative operator
before fixing \(\alpha\). Set
\begin{equation}\label{eq:amp-hodge-nucleos-relativos}
 \mathscr K^{\mathrm S}_{N+1}
 :=\ker\rho_{N+1,N},
 \qquad
 \mathscr K^{\mathrm H}_{N+1}
 :=\ker q_{N+1,N},
\end{equation}
and let \(\Delta J_{N+1}\) be the ordered family of oriented cells incorporated
when passing from \(J_N\) to \(J_{N+1}\). For a cell
\(\nu=(Z,W,a,b)\), denote its perfect
realization by \(\mathcal C_\nu=\operatorname{Cone}(\kappa_{a,b})\).
The relative cellular module is
\begin{equation}\label{eq:amp-hodge-modulo-celular-relativo}
 \mathscr A_{N+1}^{\mathrm{rel}}
 :=H^0\!\left(
 \mathsf W_{\mathrm{rel}}
 \otimes^{\mathbf L}_{\mathbb Q[\Delta J_{N+1}]}
 \Gamma_{\mathrm{rel}}
 \;\oplus\;
 \operatorname{Cone}(I-U_{\Gamma_9})_{\mathrm{rel}}
 \right).
\end{equation}
In this expression, \(\Gamma_{\mathrm{rel}}(\nu)=\mathcal C_\nu\); the
coend imposes the Mayer--Vietoris and change-of-chart relations, and the
second summand preserves the terminal difference. The relation
\(\kappa_{b,a}\tau=-P\kappa_{a,b}\) fixes the signs of the opposite
generators, while \eqref{eq:hodge-suc-carry} fixes their composition.

There are two maps, both determined by the cellular generators,
\begin{equation}\label{eq:amp-hodge-mapas-relativos}
 \begin{aligned}
  a_{N+1}&:\mathscr A_{N+1}^{\mathrm{rel}}
    \longrightarrow\mathscr K^{\mathrm S}_{N+1},\\
  b_{N+1}&:\mathscr A_{N+1}^{\mathrm{rel}}
    \longrightarrow\mathscr K^{\mathrm H}_{N+1}.
 \end{aligned}
\end{equation}
The first inserts the cone as a state correction supported on the new
charts; the second publishes its incidence evaluation. By the local
compatibility of the localized character with Mayer--Vietoris, we have
\begin{equation}\label{eq:amp-hodge-factorizacion-relativa}
 e_{N+1}^{\mathrm{rel}}a_{N+1}=b_{N+1},
 \qquad
 e_{N+1}^{\mathrm{rel}}
 :=e_{N+1}^{\mathrm{Hdg}}\big|_{\mathscr K^{\mathrm S}_{N+1}}.
\end{equation}

The relative presentation makes it possible to isolate exactly the comparison whose
exactness is equivalent to the internal surjectivity of this redundant comparator.
It is neither a hypothesis nor a source of the proprietary surjectivity already
proved in \cref{prop:amp-hodge-extension-finita}. If
\(u:\nu\to\nu'\) runs through the arrows of
\(\Delta J_{N+1}\), set
\begin{equation}\label{eq:amp-hodge-relacion-coend}
 \mathscr T_{N+1}^{\mathrm{rel}}
 :=H^0\!\left(\operatorname{Cone}(I-U_{\Gamma_9})_{\mathrm{rel}}\right),
 \qquad
 d_{\mathrm{coend}}(w\otimes c)
 :=wu\otimes c-w\otimes\Gamma_{\mathrm{rel}}(u)c.
\end{equation}
The definition of the coend provides the modules
\begin{align}
 \mathscr R_{N+1}^{\mathrm{rel}}
 &:=
 \bigoplus_{u:\nu\to\nu'}
 \mathsf W_{\mathrm{rel}}(\nu')\otimes H^0(\mathcal C_\nu),
 \label{eq:amp-hodge-modulo-relaciones}\\
 \mathscr P_{N+1}^{\mathrm{rel}}
 &:=
 \left(
  \bigoplus_{\nu\in\Delta J_{N+1}}
  \mathsf W_{\mathrm{rel}}(\nu)\otimes H^0(\mathcal C_\nu)
 \right)\oplus\mathscr T_{N+1}^{\mathrm{rel}},
 \label{eq:amp-hodge-modulo-presentacion}
\end{align}
Before comparing it with the restriction kernel, the independent relative
codomain is formed
\begin{equation}\label{eq:amp-hodge-codominio-relativo-independiente}
 \mathscr O_{N+1}^{\mathrm{rel}}
 :=\operatorname{coker}
 \left(
  d_{\mathrm{coend}}:
  \mathscr R_{N+1}^{\mathrm{rel}}
  \longrightarrow\mathscr P_{N+1}^{\mathrm{rel}}
 \right),
 \qquad
 \pi_{N+1}^{\mathrm{rel}}:
 \mathscr P_{N+1}^{\mathrm{rel}}
 \twoheadrightarrow\mathscr O_{N+1}^{\mathrm{rel}}.
\end{equation}
This definition uses only the new cells, the Mayer--Vietoris and
change-of-chart relations and the terminal; it does not use
\(\mathscr K^{\mathrm H}_{N+1}\). Let
\(\widetilde b_{N+1}:\mathscr P_{N+1}^{\mathrm{rel}}\to
\mathscr K^{\mathrm H}_{N+1}\) be the map that publishes each cellular generator
in its new incidence and the terminal generator in the terminal difference.
Since \(\widetilde b_{N+1}d_{\mathrm{coend}}=0\), there is a unique map
\begin{equation}\label{eq:amp-hodge-comparacion-codominio-relativo}
 \chi_{N+1}^{\mathrm{rel}}:
 \mathscr O_{N+1}^{\mathrm{rel}}
 \longrightarrow\mathscr K^{\mathrm H}_{N+1},
 \qquad
 \chi_{N+1}^{\mathrm{rel}}\pi_{N+1}^{\mathrm{rel}}
 =\widetilde b_{N+1}.
\end{equation}

\begin{lemma}[Canonical comparison of the relative codomain]
\label{lem:amp-hodge-comparacion-codominio-relativo}
The independent cokernel has the exact presentation
\begin{equation}\label{eq:amp-hodge-presentacion-coend-relativa}
 \mathscr R_{N+1}^{\mathrm{rel}}
 \xrightarrow{\ d_{\mathrm{coend}}\ }
 \mathscr P_{N+1}^{\mathrm{rel}}
 \xrightarrow{\ \pi_{N+1}^{\mathrm{rel}}\ }
 \mathscr O_{N+1}^{\mathrm{rel}}\longrightarrow0 .
\end{equation}
If
\(\vartheta_{N+1}:\mathscr A_{N+1}^{\mathrm{rel}}
\xrightarrow{\sim}\mathscr O_{N+1}^{\mathrm{rel}}\)
is the canonical identification of the realization \(H^0\) of the coend with its
presentation, then
\begin{equation}\label{eq:amp-hodge-b-factor-cokernel}
 b_{N+1}
 =\chi_{N+1}^{\mathrm{rel}}\vartheta_{N+1}.
\end{equation}
The comparison with all new observations is measured by
\begin{equation}\label{eq:amp-hodge-obstrucciones-comparacion}
 \mathfrak o_{N+1}^{\ker}
 :=\ker\chi_{N+1}^{\mathrm{rel}},
 \qquad
 \mathfrak o_{N+1}^{\mathrm{coker}}
 :=\operatorname{coker}\chi_{N+1}^{\mathrm{rel}}.
\end{equation}
We have
\[
 \chi_{N+1}^{\mathrm{rel}}\text{ isomorphism}
 \quad\Longleftrightarrow\quad
 \mathfrak o_{N+1}^{\ker}
 =\mathfrak o_{N+1}^{\mathrm{coker}}=0.
\]
\end{lemma}

\begin{proof}
The relations
\(wu\otimes c-w\otimes\Gamma_{\mathrm{rel}}(u)c\) publish zero because
both terms are the same incidence written in two charts. Therefore
\eqref{eq:amp-hodge-comparacion-codominio-relativo} is well defined.
The exactness of \eqref{eq:amp-hodge-presentacion-coend-relativa} is the
universal property of the cokernel. The realization \(H^0\) of the coend
provides \(\vartheta_{N+1}\), and uniqueness of factorization through the
cokernel gives \eqref{eq:amp-hodge-b-factor-cokernel}. The last equivalence
is the elementary characterization of an isomorphism by the vanishing of its
kernel and cokernel. None of these steps proves that vanishing.
\end{proof}

Thus, the coend and the restriction kernel are constructed separately. In
this auxiliary chart, relative exactness is obtained if the specific
assertion
\(\mathfrak o_{N+1}^{\ker}
=\mathfrak o_{N+1}^{\mathrm{coker}}=0\) is proved. This obligation belongs only to
this comparator: it does not govern the proprietary
APP--Mayer--Vietoris extension or the coinductive completeness already proved.

\begin{lemma}[Cellular lifting under exact comparison]
\label{lem:amp-hodge-sobreyectividad-relativa}
Suppose
\(\mathfrak o_{N+1}^{\ker}
=\mathfrak o_{N+1}^{\mathrm{coker}}=0\).
The map \(b_{N+1}\) of
\eqref{eq:amp-hodge-mapas-relativos} is surjective. Moreover, the based
order of the charts determines a right inverse
\begin{equation}\label{eq:amp-hodge-sigma-relativa}
 \sigma_{N+1}^{\mathrm{rel}}:
 \mathscr K^{\mathrm H}_{N+1}\longrightarrow
 \mathscr A_{N+1}^{\mathrm{rel}},
 \qquad
 b_{N+1}\sigma_{N+1}^{\mathrm{rel}}=I.
\end{equation}
Consequently,
\begin{equation}\label{eq:amp-hodge-seccion-relativa}
 s_{N+1}^{\mathrm{rel}}
 :=a_{N+1}\sigma_{N+1}^{\mathrm{rel}}:
 \mathscr K^{\mathrm H}_{N+1}\longrightarrow
 \mathscr K^{\mathrm S}_{N+1}
 \quad\text{satisfies}\quad
 e_{N+1}^{\mathrm{rel}}s_{N+1}^{\mathrm{rel}}=I.
\end{equation}
\end{lemma}

\begin{proof}
An element of \(\mathscr K^{\mathrm H}_{N+1}\) vanishes when restricted to
\(J_N\); therefore, it is supported on the incidences of
\(\Delta J_{N+1}\). By the hypothesis on
\eqref{eq:amp-hodge-obstrucciones-comparacion},
\(\chi_{N+1}^{\mathrm{rel}}\) is an isomorphism. The exactness of
\eqref{eq:amp-hodge-presentacion-coend-relativa} then implies that its
restriction to each new
chart is a rational combination of the classes of the cones
\(\mathcal C_\nu\). On an intersection of charts, two such expressions
differ by the element
\(wu\otimes c-w\otimes\Gamma_{\mathrm{rel}}(u)c\) of
\eqref{eq:amp-hodge-relacion-coend}; passing to the coend makes that difference
zero. If the last chart completes a nonadic turn, the difference is not
eliminated: it remains in \(\mathscr T_{N+1}^{\mathrm{rel}}\). In this way every
relative class has a preimage under \(b_{N+1}\), proving
surjectivity.

To make the right inverse explicit, the generators are ordered by
depth, support, sheet and orientation, which are data of the atlas. Row
reduction of the matrix of \(b_{N+1}\) retains the first admissible generator
of each pivot column. If
\(\{\bar c_\mu\}_\mu\) is the resulting basis of
\(\mathscr K^{\mathrm H}_{N+1}\) and \(c_\mu\) the corresponding pivot
generator, one defines
\begin{equation}\label{eq:amp-hodge-sigma-coordenada}
 \sigma_{N+1}^{\mathrm{rel}}
 \left(\sum_\mu t_\mu\bar c_\mu\right)
 :=\sum_\mu t_\mu c_\mu.
\end{equation}
Then \(b_{N+1}(c_\mu)=\bar c_\mu\), from which one obtains
\eqref{eq:amp-hodge-sigma-relativa}.

The ordering used is the restriction of the global ordering of the atlas.
If \(r:J_{N+1}\to J'_{N+1}\) is a based refinement, let
\(r_{\mathscr A}\), \(r_{\mathrm H}\) and \(r_{\mathrm S}\) be the maps
induced on the cellular module, observations and relative states.
Refinement replaces a generator with its Mayer--Vietoris
expression and does not change the first global pivot of its class. Uniqueness
of the ordered normal form gives
\begin{equation}\label{eq:amp-hodge-naturalidad-seccion-relativa}
 r_{\mathscr A}\sigma_{N+1}^{\mathrm{rel}}
 =\sigma_{N+1}^{\prime\,\mathrm{rel}}r_{\mathrm H},
 \qquad
 r_{\mathrm S}s_{N+1}^{\mathrm{rel}}
 =s_{N+1}^{\prime\,\mathrm{rel}}r_{\mathrm H}.
\end{equation}
Thus, the right inverse is compatible with refinements, not merely with an
isolated presentation. The construction is carried out in oriented pairs;
therefore it respects
\(\kappa_{b,a}\tau=-P\kappa_{a,b}\). The coend relations impose the
carry cocycle, and the terminal generator is preserved instead of being reduced to
its visible phase. Finally, \eqref{eq:amp-hodge-factorizacion-relativa}
provides \eqref{eq:amp-hodge-seccion-relativa}. No step depends on
\(\alpha\) or on an algebraic class chosen beforehand.
\end{proof}

At depth zero, the based incidences simultaneously form the bases
of \(\mathscr S_0(X,p)\) and \(\mathscr H_0^p(X)\). If
\(\widehat c_\mu\) is the seed state publishing the observation
\(c_\mu\), the map
\begin{equation}\label{eq:amp-hodge-seccion-base}
 s_0^{\mathrm{base}}
 \left(\sum_\mu t_\mu c_\mu\right)
 :=\sum_\mu t_\mu\widehat c_\mu
 \quad\text{satisfies}\quad
 e_0^{\mathrm{Hdg}}s_0^{\mathrm{base}}=I.
\end{equation}
Likewise, the unitary degeneration of the atlas defines the prolongation with
zero relative coefficients
\begin{equation}\label{eq:amp-hodge-extension-cero}
 \iota_{N+1,N}:\mathscr S_N(X,p)\longrightarrow
 \mathscr S_{N+1}(X,p),
 \qquad
 \rho_{N+1,N}\iota_{N+1,N}=I.
\end{equation}
This prolongation inherits the terminal of \(s_N\); it does not reinitialize the history.

\subsection{Auxiliary commutativity check of the relative chart}
\label{subsec:amp-hodge-control-conmutatividad-relativa}

\noindent\textbf{Status: NON\_GOVERNING\_CONTROL.}
The computable defect of the relative chart is the transformation
\begin{equation}\label{eq:amp-hodge-defecto-puente-algebraico}
 \mathfrak d_N^{\mathrm{alg}}
 :=\operatorname{cl}_{X,p}\circ\operatorname{ch}^{\mathrm{loc}}_p
   \circ R_{\mathrm{Hdg},N}
   -\operatorname{pub}_N\circ e_N^{\mathrm{Hdg}}.
\end{equation}
The preceding proprietary construction already establishes that the diagram
\eqref{eq:amp-hodge-cuadrado-finito} commutes and, therefore,
\(\mathfrak d_N^{\mathrm{alg}}=0\). The preceding expression is a computable
falsifier and a redundant check of that identity, not a hypothesis
added to the theorem.

At the increment \(N\to N+1\), let
\[
 R_{\mathrm{Hdg},N+1}^{\mathrm{rel}}
 :=R_{\mathrm{Hdg},N+1}\big|_{\mathscr K^{\mathrm S}_{N+1}},
 \qquad
 \operatorname{pub}_{N+1}^{\mathrm{rel}}
 :=\operatorname{pub}_{N+1}\big|_{\mathscr K^{\mathrm H}_{N+1}}.
\]
The two maps with the codomains required by the localized character and
by the cycle class are
\begin{equation}\label{eq:amp-hodge-mapas-relativos-tipados}
 a_{N+1}^{\mathrm{mot}}
 :=R_{\mathrm{Hdg},N+1}^{\mathrm{rel}}a_{N+1},
 \qquad
 b_{N+1}^{\mathrm{coh}}
 :=\operatorname{pub}_{N+1}^{\mathrm{rel}}b_{N+1}.
\end{equation}
The well-typed relative defect is
\begin{equation}\label{eq:amp-hodge-cuadrado-relativo-tipado}
 \mathfrak d_{N+1}^{\mathrm{rel}}
 :=
 \operatorname{cl}_{X,p}\circ\operatorname{ch}^{\mathrm{loc}}_p
 \circ a_{N+1}^{\mathrm{mot}}
 -b_{N+1}^{\mathrm{coh}}.
\end{equation}
The relative publication identity checked by this auxiliary presentation is
\begin{equation}\label{eq:amp-hodge-identidad-algebraica-requerida}
 \boxed{\mathfrak d_{N+1}^{\mathrm{rel}}=0.}
\end{equation}
It is verified on the generators that
\(R_{\mathrm{Hdg},N+1}^{\mathrm{rel}}a_{N+1}\) realizes the corresponding perfect
cone, that the localized character publishes the same incidence
as \(b_{N+1}\), that both maps annihilate the Mayer--Vietoris
relations and preserve the same terminal difference. These
equalities certify the relative chart and do not replace, condition or
reopen the governing identity already established.

\subsection{Specialized validators of the Hodge chain}

\noindent\textbf{Status: NON\_GOVERNING\_CONTROL.}
The specialized checks of that chain include the models
\(C_3/K3\), the tritic curves and the Fermat sector through the \(405\)
planar incidences and their framework of \(486\) states; the Schur--Brauer
decomposition; semistable log-Perf descent; and, for every smooth
four-dimensional cubic \(X\), the universal line correspondence
\(P\subset F(X)\times X\). If
\(p:P\to X\), \(q:P\to F(X)\) and \(g\) is the polarization, the readers
\begin{equation}\label{eq:hodge-suc-fano}
 \Phi_X=p_*q^*,
 \qquad
 \Psi_X=q_*\bigl(p^*g^2\smile-\bigr),
 \qquad
 \Psi_X\Phi_X=-6I
\end{equation}
in the corresponding sector provide an explicit rational section;
the numerators are first realized in \(\operatorname{Perf}\).

In the Weil sector, for
\(A_m=E_i^m\times E_i^m\), take
\begin{equation}\label{eq:hodge-suc-weil-graficas}
 C^1=[\Gamma_1]-[X]-[Y],
 \qquad
 C^i=[\Gamma_i]-[X]-[Y],
 \qquad
 \operatorname{cl}(C^i+iC^1)=z\wedge\bar w,
\end{equation}
and
\begin{equation}\label{eq:hodge-suc-weil-determinante}
 P_m=\det_*(C^i+iC^1),
 \qquad
 \operatorname{cl}(P_m)
 =(-1)^{m(m-1)/2}m!\,w_+.
\end{equation}
The real and imaginary parts, normalized after constructing the
perfect numerators, realize the identity in the Weil sector. For
\(m=2\), the carry is two and does not require inverting three. The
Markman--Schoen result for complex abelian varieties of dimension four and
Weil type is recorded
as a recovered and typed external result; it is not attributed as an
original construction of this work.

\noindent\textbf{Hierarchy lock.}
The three control blocks of this section take as input identities
already proved by the owner. None of their comparators, vanishings or
specialized models is a premise of
\cref{thm:amp-hodge-completitud-coinductiva,thm:hodge-suc-generacion}; there is
no return causal edge from these controls to the owner.
